\chapter{Introduction}
\label{ch:intro}

This report documents the work I carried out during my FOSSEE semester-long internship
on the \textbf{OsdagBridge} project. Before describing the technical contributions, this
chapter introduces the institutional context of the internship---the National Mission on
Education through ICT, the FOSSEE project, the Osdag software, and finally the OsdagBridge
sub-project on which all of my work was done.

\section{National Mission in Education through ICT}
The National Mission on Education through ICT (NMEICT) is a scheme under the Department of
Higher Education, Ministry of Education, Government of India. It aims to leverage the
potential of ICT to enhance teaching and learning in Higher Education Institutions in an
anytime--anywhere mode. The mission aligns with the three cardinal principles of the
Education Policy---access, equity, and quality---by:
\begin{itemize}[leftmargin=1.4em]
  \item Providing connectivity and affordable access devices for learners and institutions.
  \item Generating high-quality e-content free of cost.
\end{itemize}
NMEICT seeks to bridge the digital divide by empowering learners and teachers in urban and
rural areas, fostering inclusivity in the knowledge economy. Key focus areas include the
development of e-learning pedagogies and virtual laboratories, online testing and
certification, and the training of teachers to adopt ICT-based methods. For further details,
visit \url{www.nmeict.ac.in}.

\subsection{ICT Initiatives of MoE}
The Ministry of Education has launched several ICT initiatives aimed at students,
researchers, and institutions---such as SWAYAM, the National Digital Library, e-Yantra,
Spoken Tutorial, Virtual Labs, and FOSSEE---covering audio-video e-content, digital content
access, hands-on learning, and e-governance.

\section{FOSSEE Project}
The FOSSEE (Free/Libre and Open Source Software for Education) project promotes the use of
FLOSS tools in academia and research. It is part of the National Mission on Education through
ICT, Ministry of Education, Government of India. FOSSEE supports a range of activities such
as the Textbook Companion, Lab Migration, and the development and maintenance of open-source
scientific and engineering software.

\subsection{Internships and Fellowships}
Students from any degree and academic stage can apply for FOSSEE internships. Selection is
based on the completion of screening tasks involving programming, scientific computing, or
data collection that benefit the FLOSS community. My own selection followed the completion of
such a screening task on the Osdag codebase, after which I joined the OsdagBridge development
effort.

%% IMAGE NEEDED: fossee_activities.png -- (optional) FOSSEE projects & activities graphic
\osdagfig{fossee_activities.png}{FOSSEE projects and activities.}{fig:fossee}

\section{Osdag Software}
Osdag (Open Steel Design and Graphics) is a cross-platform, free/libre and open-source
software for the detailing and design of steel structures based on the Indian Standard
IS~800:2007. It allows users to design steel connections, members, and systems through an
interactive graphical user interface (GUI) and provides 3D visualizations of the designed
components. The software enables easy export of CAD models to drafting tools for
construction and fabrication drawings. Built on Python and several Python-based FLOSS tools
(notably PyQt and PythonOCC), Osdag is licensed under the GNU Lesser General Public License
(LGPL) Version~3.

\subsection{Osdag GUI}
The Osdag GUI is designed to be user-friendly and interactive. It consists of:
\begin{itemize}[leftmargin=1.4em]
  \item \textbf{Input Dock:} Collects and validates user inputs.
  \item \textbf{Output Dock:} Displays design results after validation.
  \item \textbf{CAD Window:} Displays the 3D CAD model, where users can pan, zoom, and rotate
        the design.
  \item \textbf{Message Log:} Shows errors, warnings, and suggestions based on design checks.
\end{itemize}
Each of these four areas features directly in my internship work: I extended the CAD Window
with node and grillage visualisation and a full navigation toolbar, added validators to the
Input Dock, contributed to the design-report generation pipeline, and added a dedicated log
chapter to the report.

\section{OsdagBridge}
\textbf{OsdagBridge} is the bridge-design extension of the Osdag ecosystem. Where the core
Osdag software targets steel connections and members, OsdagBridge focuses on the design of
steel--concrete composite \emph{plate-girder bridges}. It reuses the Osdag GUI paradigm---an
Input Dock, an Output Dock, a 3D CAD window, and a message log---while adding
bridge-specific geometry (girders, cross bracings, end diaphragms, deck slab, shear studs,
railings, and medians), plate-girder design calculations, and an automatically generated,
IS-code-based design report.

All of the work described in this report was contributed to the OsdagBridge repository
maintained at \texttt{Aditya-Donde/OsdagBridge} under the username \texttt{OmPathania21}.

\section{Scope of this Internship}
My contributions span the full breadth of the OsdagBridge application and are grouped, in the
order in which I worked on them, as follows:
\begin{enumerate}[leftmargin=1.6em]
  \item \textbf{3D CAD viewer} --- node information on hover, node numbering, a grillage view
        connecting the structural nodes, and show/hide toggles (Chapter~\ref{ch:cad}).
  \item \textbf{Toolbar system} --- a dedicated toolbar controller wiring zoom, pan, rotate,
        and zoom-window camera operations, together with label-visibility toggles and later
        state-management fixes (Chapter~\ref{ch:toolbar}).
  \item \textbf{Bridge geometry} --- dynamic geometry keys and a shear-stud solver in the
        plate-girder bridge module (Chapter~\ref{ch:geometry}).
  \item \textbf{Coordinate axis} --- an always-visible coordinate axis in the 3D view with
        context-aware toolbar filtering (Chapter~\ref{ch:axis}).
  \item \textbf{Design report generation} --- the single largest part of my work: building
        multiple chapters, sections, and summary tables of the automatically generated PDF
        design report, along with report customisation and a log-window chapter
        (Chapter~\ref{ch:report}).
  \item \textbf{Report plot capture and zoom} --- capturing Matplotlib plots into the report
        and a zoom window for plots (Chapter~\ref{ch:plots}).
  \item \textbf{UI/UX enhancements} --- a custom cursor, a legend toggle, and conditional
        visibility for optional bridge components (Chapter~\ref{ch:ui}).
  \item \textbf{Material input validators} --- validation of custom material properties in the
        Input Dock (Chapter~\ref{ch:validators}).
\end{enumerate}

The remainder of this report describes each of these areas in turn, with the relevant code
from my pull requests, followed by a concluding chapter of lessons learned.
